% Zufallsexperimente und Pfadregeln – bank/zufallsexperimente-und-pfadregeln/ (zusammenbau v0.4, Heft)
\einheitenkopf[t1]{Zufallsexperimente und Pfadregeln}
\setcounter{aufgabe}{0}

% Anlauf (ohne Prüfkennung)
\begin{aufgabe}{Ereignisse als Mengen – Anlauf}
\begin{teile}
\teil Zwei Würfe mit einem Würfel: „Erst eine 2, dann eine 6.“ Ergebnis oder Ereignis?\\ \kreuz{einzelnes Ergebnis}\\ \kreuz{Ereignis aus mehreren Ergebnissen}
\teil Eine Münze wird dreimal geworfen (K = Kopf, Z = Zahl). Ergebnismenge als Tupel? Laplace-Experiment?
\teil Eine Münze wird dreimal geworfen. A: „mindestens zweimal Kopf“, B: „der erste Wurf ist Zahl“. Ergebnisse von $A \cap B$ und von „weder A noch B“?
\end{teile}
\end{aufgabe}

\begin{aufgabe}{Ereignisse als Mengen – Prüfungsaufgaben}
\begin{teile}
\teil Fahrradcheck: 25\,\% der Räder haben einen Mangel am Licht (L), 15\,\% an der Bremse (B), 68\,\% haben keinen der beiden Mängel. Weise nach, dass 8\,\% beide Mängel haben. \hfill (Abitur 2023 GK)
\teil Handyreparatur: 20\,\% der Geräte haben ein defektes Display (D), 30\,\% einen schwachen Akku (A), 60\,\% keinen der beiden Mängel. Weise nach, dass 10\,\% beide Mängel haben. \hfill (Abitur 2023 GK)
\end{teile}
\end{aufgabe}

% Anlauf (ohne Prüfkennung)
\begin{aufgabe}{Laplace-Experimente – Anlauf}
\begin{teile}
\teil Ein Glücksrad hat vier gleich große Sektoren. Sind die vier Ergebnisse gleich wahrscheinlich?\\ \kreuz{ja, Laplace}\\ \kreuz{nein, kein Laplace}
\teil In einer Lostrommel liegen 50 Lose, 8 davon sind Gewinne. P(Gewinn)? P = \leerfeld
\teil Aus einem Beutel mit 30 gleich großen, gleich schweren Murmeln wird blind eine gezogen. Begründe, dass das ein Laplace-Experiment ist, und nenne ein weiteres Beispiel.
\end{teile}
\end{aufgabe}

\begin{aufgabe}{Laplace-Experimente – Prüfungsaufgaben}
\begin{teile}
\teil Eine Urne enthält n Kugeln, 20\,\% davon tragen einen Stern. Zwei Kugeln mit Stern werden hinzugelegt, dann wird eine Kugel gezogen. Weise über die Differenz der Wahrscheinlichkeiten nach, dass P(Stern) größer geworden ist. \hfill (Abitur 2019 GK)
\teil Eine Tüte enthält n Bonbons, 30\,\% davon schmecken nach Zitrone. Ein Zitronen- und zwei Kirschbonbons kommen dazu, dann wird eines gezogen. Weise über die Differenz nach, dass P(Zitrone) zunimmt. \hfill (Abitur 2019 GK)
\teil Eine Urne enthält n Kugeln, 40\,\% davon blau. Eine blaue und zwei rote Kugeln werden hinzugelegt. Begründe, dass $\frac{0{,}4n + 1}{n + 3}$ die Wahrscheinlichkeit für eine blaue Kugel angibt. \hfill (Abitur 2019 GK)
\teil Ein Glas enthält n Murmeln, ein Viertel davon golden. Zwei goldene und zwei silberne kommen dazu. Begründe, dass $\frac{0{,}25n + 2}{n + 4}$ die Wahrscheinlichkeit für eine goldene Murmel angibt. \hfill (Abitur 2019 GK)
\end{teile}
\end{aufgabe}

% Anlauf (ohne Prüfkennung)
\begin{aufgabe}{Pfadregeln bei unabhängigen Stufen – Anlauf}
\begin{teile}
\teil Ein Glücksrad wird dreimal gedreht. Bleibt P(rot) auf jeder Stufe gleich?\\ \kreuz{p bleibt gleich}\\ \kreuz{p ändert sich}
\teil Ein Glücksrad zeigt rot (R) mit 30\,\%, sonst blau (B); es wird zweimal gedreht. Beschrifte den Baum. P(erst rot, dann blau)?

\baumzwei{R/,B/}{R/,B/}{R/,B/}
\teil Ein Glücksrad liefert einen Treffer mit $0{,}4$. Es wird fünfmal gedreht. P(erst drei Nieten, dann zwei Treffer)? P = \leerfeld
\end{teile}
\end{aufgabe}

\begin{aufgabe}{Pfadregeln bei unabhängigen Stufen – Prüfungsaufgaben}
\begin{teile}
\teil 15\,\% aller Bienenstöcke sind von Milben befallen. Zwei Imker kontrollieren je 8 zufällig gewählte Stöcke. P(mindestens einer der beiden findet keinen befallenen Stock)? \hfill (Abitur 2022 GK) \\ P = \leerfeld
\teil 20\,\% der Hausaufgabenhefte sind unvollständig. Zwei Lehrkräfte prüfen je 6 zufällig gewählte Hefte. P(mindestens eine der beiden findet nur vollständige Hefte)? \hfill (Abitur 2022 GK) \\ P = \leerfeld
\teil Würfel Grün trägt 1, 2, 3, 4, 4, 6, Würfel Rot trägt 1, 2, 3, 4, 6, 6. Mit jedem wird zweimal geworfen. Rot behauptet, Augensumme 8 sei mit Rot wahrscheinlicher. Widerlege über die beitragenden Pfade. \hfill (Abitur 2020 GK)
\teil Würfel Blau trägt 1, 1, 2, 3, 5, 6, Würfel Gelb trägt 1, 2, 3, 5, 5, 6. Mit jedem wird zweimal geworfen. Gelb behauptet, Augensumme 7 sei mit Gelb wahrscheinlicher. Widerlege über die beitragenden Pfade. \hfill (Abitur 2020 GK)
\end{teile}
\end{aufgabe}

% Anlauf (ohne Prüfkennung)
\begin{aufgabe}{Ohne Zurücklegen und Umlegen – Anlauf}
\begin{teile}
\teil Aus einer Klasse werden nacheinander zwei verschiedene Kinder für zwei Ämter ausgelost. Bleibt P(Mädchen) beim zweiten Losen gleich?\\ \kreuz{p bleibt gleich}\\ \kreuz{p ändert sich}
\teil In einer Klasse mit 25 Kindern tragen 10 eine Brille. Zwei Kinder werden ausgelost. P(beide mit Brille)? P = \leerfeld
\teil Aus einem Skatblatt (32 Karten, 8 Herz) werden vier Karten gezogen. P(alle vier Herz)? P = \leerfeld
\end{teile}
\end{aufgabe}

\begin{aufgabe}{Ohne Zurücklegen und Umlegen – Prüfungsaufgaben}
\begin{teile}
\teil In einem Hut liegen 7 Lose, 3 davon Freikarten. Sieben Kinder ziehen nacheinander ohne Zurücklegen. Hat das erste oder das dritte Kind die größere Chance auf eine Freikarte? Begründe. \hfill (Abitur 2021 GK)
\teil Unter 12 Umschlägen stecken 4 Gutscheine. Eva zieht als Zweite, Tom als Fünfter, ohne Zurücklegen und ohne die früheren Ergebnisse zu kennen. Wer hat die bessere Chance? Begründe. \hfill (Abitur 2021 GK)
\end{teile}
\end{aufgabe}

% Anlauf (ohne Prüfkennung)
\begin{aufgabe}{Rückwärts – Anlauf}
\begin{teile}
\teil „Eine Urne enthält 3 rote und 5 blaue Kugeln. Wie wahrscheinlich sind zwei rote beim Ziehen ohne Zurücklegen?“\\ \kreuz{vorwärts: eine Wahrscheinlichkeit ist gesucht}\\ \kreuz{rückwärts: eine Größe im Baum ist gesucht}
\teil 60\,\% der Kunden sind Frauen, unter ihnen sind 50\,\% Stammkunden. Insgesamt sind 44\,\% Stammkunden. Anteil a der Stammkunden unter den Männern? a = \leerfeld
\teil Befragung mit Zufallsfrage: Mit $0{,}7$ kommt F1 „Hast du schon einmal …?“, mit $0{,}3$ F2 „Hast du noch nie …?“. 50\,\% antworten mit Ja. Anteil p derer, die es schon getan haben? p = \leerfeld
\end{teile}
\end{aufgabe}

\begin{aufgabe}{Rückwärts – Prüfungsaufgaben}
\begin{teile}
\teil Ein Glücksrad hat einen grünen Sektor mit P(grün) $= 0{,}4$, dazu einen roten und einen gelben. Es wird zweimal gedreht. A: einmal Rot und einmal Gelb. Der rote Sektor ist so gewählt, dass P(A) maximal ist. Mittelpunktswinkel des roten Sektors? \hfill (Abitur 2022 GK) \\ \leerfeld[°]
\teil Ein Glücksrad hat einen weißen Sektor mit P(weiß) $= \frac{1}{3}$, dazu einen schwarzen und einen grauen. Es wird zweimal gedreht. A: einmal Schwarz und einmal Grau. Der schwarze Sektor ist so gewählt, dass P(A) maximal ist. Mittelpunktswinkel des schwarzen Sektors? \hfill (Abitur 2022 GK) \\ \leerfeld[°]
\end{teile}
\end{aufgabe}
